% math4cs-hw9-planarity-coloring-solution.tex

% !TEX program = xelatex
%%%%%%%%%%%%%%%%%%%%
% see http://mirrors.concertpass.com/tex-archive/macros/latex/contrib/tufte-latex/sample-handout.pdf
% for how to use tufte-handout
\documentclass[a4paper, justified]{tufte-handout}

\input{hw-preamble} % feel free to modify this file if you understand LaTeX well
%%%%%%%%%%%%%%%%%%%%
\title{计算机数学 (10-planarity-coloring: 平面图与图着色)}
\me{魏恒峰}{hfwei@hnu.edu.cn}{}{}
\date{2026年5月30日 发布作业 \qquad 2026年6月19日 发布答案}
%%%%%%%%%%%%%%%%%%%%
\begin{document}
\maketitle
%%%%%%%%%%%%%%%%%%%%
\noplagiarism % PLEASE DON'T DELETE THIS LINE!
%%%%%%%%%%%%%%%%%%%%
\begin{abstract}
\end{abstract}
%%%%%%%%%%%%%%%%%%%%
\beginrequired
%%%%%%%%%%%%%%%

%%%%%%%%%%%%%%%
\begin{problem}[平面图]
  假设 $G$ 是顶点数 $\ge 11$ 的简单图, $\overline{G}$ 是 $G$ 的补图~\footnote{
    补图: 顶点集相同, 但是 $e$ 是 $G$ 的边当且仅当 $e$ 不是 $\overline{G}$ 的边。
  }。

  \begin{enumerate}[(1)]
    \item 请证明: $G$ 和 $\overline{G}$ 不同为平面图。
    \item 请举例说明: 当 $G$ 的顶点数为 $8$ 时,
      $G$ 和 $\overline{G}$ 可以同时为平面图。
  \end{enumerate}
\end{problem}

\begin{proof}
  \begin{enumerate}[(1)]
    \item 反设 $G$ 和 $\overline{G}$ 都是平面图。因此,
      \begin{align}
        m \le 3n - 6 \\
        \overline{m} \le 3n - 6 \\
        m + \overline{m} = \frac{n(n-1)}{2}
      \end{align}
      解得,
      \[
        \frac{13 -\sqrt{73}}{2} \le n \le \frac{13 + \sqrt{73}}{2}.
      \]
      其整数解为
      \[
        n \in \mathbb{Z} \land 3 \le n \le 10.
      \]
      与 $n \ge 11$ 矛盾。

    \item 取顶点集 $V = \set{1,2,3,4,5,6,7,8}$, 边集 $E(G)$ 由以下 $12$ 条边组成:
      \begin{itemize}
        \item 两个 $4$-圈: $12, 23, 34, 41$ 与 $56, 67, 78, 85$
          (此处 $ij$ 表示顶点 $i$ 与 $j$ 之间的边);
        \item 连接两部分的 $4$ 条边: $15, 16, 27, 36$。
      \end{itemize}
  \end{enumerate}
\end{proof}
%%%%%%%%%%%%%%%

%%%%%%%%%%%%%%%
\begin{problem}[图同胚]
  假设图 $G_{1}$ 与 $G_{2}$ 是 homeomorphic (同胚) 的。
  请证明~\footnote{$m$, $n$ 分别表示边数与点数。}:
  \[
    m_{1} - n_{1} = m_{2} - n_{2}.
  \]
\end{problem}

\begin{proof}
  分别考虑``插入''与``收缩''2度顶点的操作:
  \begin{itemize}
    \item 对于``插入''一个2度顶点, 顶点数 $n$ 加一, 边数 $m$ 加一, 所以 $m-n$ 保持不变。
    \item 对于``收缩''一个2度顶点, 顶点数 $n$ 减一, 边数 $m$ 减一, 所以 $m-n$ 保持不变。
  \end{itemize}
  对从 $G_{1}$ 转化到 $G_{2}$ 的``插入/收缩''操作序列的长度作归纳, 即得证。
\end{proof}
%%%%%%%%%%%%%%%

%%%%%%%%%%%%%%%
\begin{problem}[图着色]
  假设 $G$ 是包含 $n$ 个顶点的$d$-正则简单图。

  \noindent 请证明
  \[
    \chi(G) \ge \frac{n}{n-d}\text{。}
  \]
\end{problem}

\begin{proof}
  反设
  \[
    \chi(G) < \frac{n}{n-d}.
  \]
  因此, 将 $G$ 中顶点按照颜色归类, 可以将 $G$ 划分成 $< \frac{n}{n-d}$ 组,
  每组内的顶点两两不相邻~\footnote{注, 这种组被称为``独立集''}。
  根据鸽笼原理, 至少存在一组含有 $> \frac{n}{\frac{n}{n-d}} = n - d$ 个顶点。
  由于 $G$ 是 $d$-正则简单图, 该组中必有两顶点相邻。矛盾。
\end{proof}
%%%%%%%%%%%%%%%

%%%%%%%%%%%%%%%
\begin{problem}[平面图的着色]
  假设 $G$ 是不包含三角形 $\triangle$ 的简单平面图。
  \begin{enumerate}[(1)]
    \item 请使用Euler公式证明 $G$ 含有度数 $\le 3$ 的顶点。
    \item 请使用数学归纳法证明 $G$ 是 $4$-可着色的。
  \end{enumerate}
\end{problem}

\begin{proof}
  \begin{enumerate}[(1)]
    \item 反设 $\delta(G) \ge 4$。因此,
      \[
        \sum_{v} \deg(v) \ge 4n.
      \]
      由于 $G$ 是不包含$\triangle$的简单平面图,
      \[
        m \le 2n - 4.
      \]
      根据握手定理,
      \[
        \sum_{v} \deg(v) = 2m \le 4n - 8.
      \]
      矛盾。
    \item 对 $G$ 中顶点数 $n$ 作归纳。
      \begin{description}
        \item[基础步骤:] $n = 1$。$G$ 显然是 $4$-可着色的。
        \item[归纳假设:] 假设命题对包含 $1 \le k = n-1$ 个顶点的任意符合要求的图均成立。
        \item[归纳步骤:] 考虑包含 $n$ 个顶点的图 $G$。根据 (1), $G$ 中包含度数 $\le 3$
          的顶点, 记为 $v$。从 $G$ 中删除 $v$, 得到图 $G'$。$G'$ 包含 $n-1$ 个顶点,
          且是不包含$\triangle$的简单平面图。根据归纳假设, $G'$ 是 $4$-可着色的。
          考虑 $G'$ 的任意一种 $4$-着色方案, 将 $v$ (及删除的边) 加入 $G'$, 重新得到图 $G$。
          因为 $\deg(v) \le 3$, 所以 $G$ 是 $4$-可着色的。
      \end{description}
  \end{enumerate}
\end{proof}
%%%%%%%%%%%%%%%

%%%%%%%%%%%%%%%
\begin{problem}[着色多项式]
  设 $C_{n}$ 是 $n$ 个顶点的圈图。
  请给出 $C_{n}$ 的着色多项式 $P(C_{n}, k)$。

  \noindent 要求给出证明或计算过程。
\end{problem}

\begin{solution}
  使用删点-缩边递推公式。设 $e$ 是 $C_n$ 的一条边, 则
  \[
    P(C_n, k) = P(C_n - e, k) - P(C_n / e, k).
  \]
  其中 $C_n - e = P_n$ 是 $n$ 个顶点的路径, $C_n / e = C_{n-1}$。
  又已知
  \[
    P(P_n, k) = k(k-1)^{n-1}.
  \]
  故
  \[
    P(C_n, k) = k(k-1)^{n-1} - P(C_{n-1}, k).
  \]
  令 $a_n = P(C_n,k)$, 则
  \[
    a_n = k(k-1)^{n-1} - a_{n-1},
    \qquad
    a_1 = k.
  \]
  解得
  \[
    P(C_n, k) = (k-1)^n + (-1)^n (k-1).
  \]
\end{solution}
%%%%%%%%%%%%%%%

%%%%%%%%%%%%%%%
\begin{problem}[图着色的应用]
  现在需要为七场讲座制定一份时间表。
  由于有些学生希望参加多场讲座，某些讲座不能安排在同一时间段。
  下表中用星号标出了不能同时进行的讲座对。

  \noindent 请问，安排这七场讲座至少需要多少个时间段？
  \fig{width = 0.50\textwidth}{figs/timetable}
\end{problem}

\begin{solution}
  将七场讲座看作顶点, 若两场讲座不能同时进行则连边, 得到冲突图 $G$。
  所需最少时间段数就是 $\chi(G)$。

  子图 $\set{a,b,c,d}$ 构成 $K_4$, 故 $\chi(G) \ge 4$。

  下面给出一个 $4$-着色 (冲突边为
  $ab, ac, ad, ag, bc, bd, be, bg, cd, cf, df, fg$):
  \[
    a \mapsto 1,\quad b \mapsto 2,\quad c \mapsto 3,\quad d \mapsto 4,
  \]
  \[
    e \mapsto 1,\quad f \mapsto 1,\quad g \mapsto 3.
  \]
  直接检查可知, 任意冲突对颜色不同。
  因而 $\chi(G) = 4$。

  因此, 至少需要 $4$ 个时间段。
\end{solution}
%%%%%%%%%%%%%%%
%%%%%%%%%%%%%%%%%%%%
\end{document}